% SGA 3，第 VII 篇_B，英文重建本——第 1 轮（中文译本）。
% 范围：依据原文件第 23--26 页，完整包含 1.2.4 至 1.2.6；止于原文件第 26 页的 1.3 之前。
% 合并阅读版页码： 536--539. 再版印刷页码： 525--528.
% 本范围内的原版页边定位： 497--498.
% 依据： Polo--Gille Exp7B-13oct24.pdf,
% SHA-256 D303690E0195654BBE7C3DAA5D89EA23E24CD709A3575217234AC19AE8836694.
% 曾参阅 Jacob Reinhold 的 CC BY 4.0 Markdown，作为对比版本来源。
% 本范围内无原生图；所有陈列数学均为原生 LaTeX。

\medskip
\noindent\SGAInlineRefTarget{sga3:vii:paragraph:1-2-4}\textbf{1.2.4.}\footnote{编者注：本段的顺序已经调整：先定义 \(\widehat S_k(E)\)，再说明
\(\mathbf V^f_k(E)\) 表示 \(\mathbf V^f_k(E)\)。}
\phantomsection
\label{sga3:VIIB:1.2.4}
对于每个伪紧 \(k\)-模 \(E\)，以 \(\widehat S_k(E)\) 表示 \(E\) 的
\emph{完备化对称代数}，其定义如下。令 \(T_k(E)\) 为下列伪紧 \(k\)-模的直和
\[
  \widehat{\bigotimes}_{k}^{\,n}E
  =
  E\widehat{\otimes}_{k}\cdots\widehat{\otimes}_{k}E
  \qquad
  \left(
    n\geq0;\quad
    \widehat{\bigotimes}_{k}^{\,0}E=k
  \right).
\]
赋予 \(T_k(E)\) 通常的分次 \(k\)-代数乘法。令 \(S_k(E)\) 为
\(T_k(E)\) 关于如下齐次理想的商：其第 \(n\) 个分量是
\(\widehat{\bigotimes}_{k}^{\,n}E\) 中由下列元素生成的闭 \(k\)-子模
\[
\begin{aligned}
&x_1\widehat{\otimes}\cdots\widehat{\otimes}x_i
 \widehat{\otimes}x_{i+1}\widehat{\otimes}\cdots
 \widehat{\otimes}x_n
\\[-0.1em]
&\qquad{}-
x_1\widehat{\otimes}\cdots\widehat{\otimes}x_{i+1}
 \widehat{\otimes}x_i\widehat{\otimes}\cdots
 \widehat{\otimes}x_n.
\end{aligned}
\]
若将这个商记为 \(S_k^n(E)\)，则 \(S_k(E)\) 是第 \(n\) 个分量为
\(S_k^n(E)\) 的分次 \(k\)-代数。

在 \(S_k(E)\) 上赋予由理想 \(\mathfrak l\) 的集合 \(\Upsilon\) 定义的线性拓扑，
其中要求 \(S_k(E)/\mathfrak l\) 是有限长度的 \(k\)-模，且对每个 \(n\)，
\[
  S_k^n(E)\cap\mathfrak l
\]
都是 \(S_k^n(E)\) 的开子模。预有限代数 \(\widehat S_k(E)\) 是
\(S_k(E)\) 对于此拓扑的分离完备化。\footnote{编者注：例如，设 \(k\) 是域，
\(E\) 是伪紧 \(k\)-向量空间，令 \(V=\operatorname{Hom}_c(E,k)\)；则
\(E\simeq V^*\)\textup{（参见\ \SGARefLink{sga3:VIIB:0.2.2}{0.2.2}）}。
对 \(V\) 的每个有限维子空间 \(W\)，令 \(\mathcal F(W)\) 为
\(k\)-概形 \(\mathbf V(W)=\operatorname{Spec}S_k(W^*)\) 的闭点集，
并令 \(\mathcal F(V)=\varinjlim_W\mathcal F(W)\)。则 \(\widehat S_k(E)\)
是以下伪紧 \(k\)-代数对 \(x\in\mathcal F(V)\) 的积
\[
  \widehat S_k(E)_x
  =
  \varprojlim_W\widehat{\mathcal O}_{W,x},
\]
其中 \(W\) 遍历 \(V\) 中满足 \(x\in\mathcal F(W)\) 的有限维子空间，
而 \(\widehat{\mathcal O}_{W,x}\) 是局部环 \(\mathcal O_{\mathbf V(W),x}\)
对于由有限余维理想定义的拓扑的分离完备化\textup{（在这里，该拓扑与
\(\mathfrak m\)-进拓扑一致）}。若 \(k\) 是代数闭域，且 \((v_i)_{i\in I}\) 是
\(V\) 的基，从而 \(E\) 有伪基 \((e_i)_{i\in I}\)，则每个局部分量
\(\widehat S_k(E)_x\) 都同构于赋予由有限余维理想定义的拓扑的形式幂级数环
\(k[[e_i,\ i\in I]]\)。}

我们以
\[
  \mathbf V^f_k(E)
  =
  \operatorname{Spf}\bigl(\widehat S_k(E)\bigr)
\]
表示相应的形式 \(k\)-簇。它表示 \(k\)-函子 \(\mathbf V^f_k(E)\)；
即对于 \(\operatorname{Alf}/k\) 的每个对象 \(C\)，有典范同构
\[
\begin{aligned}
  \mathbf V^f_k(E)(C)
  &=
  \operatorname{Hom}_c(E,C|_k)
  \xrightarrow{\ \sim\ }
  \operatorname{Hom}_{\operatorname{Alp}/k}
    \bigl(\widehat S_k(E),C\bigr)
\\
  &=
  \operatorname{Hom}_{\operatorname{Vaf}/k}
    \bigl(\operatorname{Spf}(C),\mathbf V^f_k(E)\bigr).
\end{aligned}
\]
下文将形式簇 \(\mathbf V^f_k(E)\) 与 \(k\)-函子 \(\mathbf V^f_k(E)\) 等同。

\medskip
\noindent\textbf{1.2.5.}
\phantomsection
\label{sga3:VIIB:1.2.5}
由 \SGARefLink{sga3:vii:paragraph:1-2-4}{1.2.4}，从 \(E\) 到 \(k\) 的零态射
对应于预有限代数态射
\[
  \pi:\widehat S_k(E)\longrightarrow k.
\]
态射 \(\pi\) 在 \(n\geq1\) 时在 \(S_k^n(E)\) 上诱导零映射，并定义结构态射
\[
  \mathbf V^f_k(E)\longrightarrow\operatorname{Spf}(k).
\]

的截面。我们以 \(\mathbf V^{f,0}_k(E)\) 表示如下形式簇：其点是
\(\operatorname{Spf}(k)\) 的点在截面 \(\operatorname{Spf}(\pi)\) 下的像，
这些点处的局部代数与 \(\mathbf V^f_k(E)\) 的相应局部代数相同。
\footnote{编者注：本段的其余部分已经详细展开。}

因此，\(\mathbf V^{f,0}_k(E)\) 的仿射代数是 \(S_k(E)\) 对于由下列理想
\(\mathfrak l\in\Upsilon\)\textup{（参见\ \SGARefLink{sga3:vii:paragraph:1-2-4}{1.2.4}）}
定义的拓扑的分离完备化：对所有充分大的 \(n\)，该理想包含
\(S_k^n(E)\)。因此，它是无穷积
\[
  k[[E]]
  =
  k\times E\times S_k^2(E)\times S_k^3(E)\times\cdots.
\]

设 \(C\) 是 \(\operatorname{Alf}/k\) 的对象，设
\[
  u\in\mathbf V^f_k(E)(C)=\operatorname{Hom}_c(E,C),
\]
并设
\[
  \phi:\widehat S_k(E)\longrightarrow C
\]
是相应的预有限 \(k\)-代数态射。则 \(\operatorname{Ker}(\phi)\) 是开理想，
即 \(\phi\) 属于 \(\mathbf V^{f,0}_k(E)(C)\)，当且仅当
\(\operatorname{Ker}(\phi)\) 对所有充分大的 \(n\) 都包含 \(S_k^n(E)\)，
或等价地，当且仅当 \(u(E)\) 包含于 \(C\) 的根 \(r(C)\) 中。
因此，对 \(\operatorname{Alf}/k\) 的每个对象 \(C\)，有典范同构
\[
  \mathbf V^{f,0}_k(E)(C)
  \simeq
  \operatorname{Hom}_{\operatorname{Alp}/k}(k[[E]],C)
  \simeq
  \operatorname{Hom}_c(E,r(C)).
\]

\medskip
\noindent\textbf{1.2.6.}
\phantomsection
\label{sga3:VIIB:1.2.6}
若形式 \(k\)-簇 \(V\) 的仿射代数长度有限，则称其\emph{具有有限长度}。
类似地，若 \(S\) 是概形，则称 \(S\)-概形 \(X\) \emph{具有有限长度}，如果
\(X\) 在 \(S\) 上有限，且 \(\mathcal O_X\) 在 \(S\) 上的直像是有限长度的
\(\mathcal O_S\)-模。\footnote{编者注：下句为新增。}
因此，给定一个有限长度的 \(S\)-概形 \(X\)，等价于给定 \(S\) 的一个有限闭点集
\[
  \{s_1,\ldots,s_n\}
\]
以及每个点处的一个有限长度的 \(\mathcal O_{S,s_i}\)-代数。

由此可知，有限长度的 \(S\)-概形等同于下列形式概形 \(\widehat S\) 上的
有限长度形式簇。\(\widehat S\) 的底拓扑空间是 \(S\) 的闭点集，赋予离散拓扑。
若 \(s\) 是这样的闭点，则结构层的茎 \(\mathcal O_{\widehat S,s}\) 是
\(\mathcal O_{S,s}\) 对于由有限余长度理想定义的线性拓扑的分离完备化
\[
  \widehat{\mathcal O}_{S,s}
\]
因此
\[
  \widehat S
  =
  \operatorname{Spf}\mathcal A(\widehat S),
\]
其中 \(\mathcal A(\widehat S)\) 是各 \(\widehat{\mathcal O}_{S,s}\) 的积，
\(s\) 遍历 \(S\) 的闭点，并赋予积拓扑。

\medskip
\noindent\textbf{定义。}
\phantomsection
\label{sga3:VIIB:1.2.6.definition}
\textit{若 \(X\) 是 \(S\)-概形，以 \(\widehat X/\widehat S\) 表示如下定义的
\[
  k=\mathcal A(\widehat S)
\]
上的形式簇。其底拓扑空间由满足
\[
  [\kappa(x):\kappa(s)]<\infty,
\]
的点 \(x\in X\) 组成，其中 \(s\) 是 \(x\) 在 \(S\) 中的像。
局部环 \(\mathcal O_{\widehat X/\widehat S,x}\) 是 \(\mathcal O_{X,x}\)
对于由 \(\mathcal O_{X,x}\) 的如下理想 \(I\) 定义的线性拓扑的分离完备化：
\(\mathcal O_{X,x}/I\) 作为 \(\mathcal O_{S,s}\)-模长度有限。由于
\([\kappa(x):\kappa(s)]<\infty\)，这等价于要求 \(\mathcal O_{X,x}/I\)
作为 \(\mathcal O_{X,x}\)-模长度有限。}

设 \(\operatorname{Vaf}^{\ell f}/\widehat S\) 是 \(\widehat S\) 上的有限长度形式簇范畴，
它与有限长度的 \(S\)-概形范畴等同。由 \SGARefLink{sga3:VIIB:1.1}{1.1} 和
\SGARefLink{sga3:VIIB:1.2.1}{1.2.1}，\(\widehat S\) 上的形式簇范畴
\(\operatorname{Vaf}/\widehat S\) 等价于以下左正合反变函子所成的范畴
\[
  \operatorname{Vaf}^{\ell f}/\widehat S
  \longrightarrow
  (\mathbf{Sets}).
\]
特别地，对于每个 \(S\)-概形 \(X\)，函子
\[
  T\longmapsto
  \operatorname{Hom}_{(\operatorname{Sch}/S)}(T,X)
\]
从 \(\operatorname{Vaf}^{\ell f}/\widehat S\) 到 \((\mathbf{Sets})\) 是左正合的，
因而对应于 \(\widehat S\) 上的一个形式簇。这个簇恰是
\(\widehat X/\widehat S\)。\footnote{编者注：下面的命题已作扩充，加入了函子
\(X\mapsto\widehat X/\widehat S\) 与有限投射极限交换这一事实。在原版中，
这出现在 \SGARefLink{sga3:VIIB:1.3.4}{1.3.4} 的证明中；这里给出的证明比原证明更直接。
这一结果还将在\SGARefLink{sga3:VIIB:2}{第~2 节}和
\SGARefLink{sga3:VIIA:3.3.1}{3.3.1} 中用到。}

\medskip
\noindent\textbf{命题。}
\phantomsection
\label{sga3:VIIB:1.2.6.proposition}
\textit{对于每个 \(S\)-概形 \(X\)，函子}
\[
  \operatorname{Hom}_{\operatorname{Vaf}/\widehat S}
    (-,\widehat X/\widehat S)
  \qquad\textit{和}\qquad
  \operatorname{Hom}_{(\operatorname{Sch}/S)}(-,X)
\]
\textit{在 \(\operatorname{Vaf}^{\ell f}/\widehat S\) 上的限制相同。这定义了从
\((\operatorname{Sch}/S)\) 到 \(\operatorname{Vaf}/\widehat S\) 的函子}
\[
  X\longmapsto\widehat X/\widehat S
\]
\textit{而且这个函子与有限投射极限交换。}

\noindent\textit{证明。}
形式簇 \(\widehat X/\widehat S\) 具有所要求的性质，且对应
\(X\mapsto\widehat X/\widehat S\) 是函子性的，这都是直接的。只需证明第二个断言。

令
\[
  \mathscr S=(\operatorname{Sch}/S),
  \qquad
  \mathscr V=\operatorname{Vaf}/\widehat S,
\]
并将 \(\widehat X/\widehat S\) 简记为 \(\widehat X\)。
由 \SGARefLink{sga3:VIIB:1.2.B}{1.2.B}，范畴 \(\mathscr V\) 有任意投射极限。
设 \((X_i)_{i\in I}\) 是 \(\mathscr S\) 中的投射系，并假定
\[
  X=\varprojlim_iX_i
\]
在 \(\mathscr S\) 中存在，\(I\) 有限时确实如此。
将 \(\mathscr S\) 的每个对象 \(Y\) 对应于函子
\[
  h_Y=\operatorname{Hom}_{\mathscr S}(-,Y)
\]
以及将 \(\mathscr V\) 的每个对象 \(V\) 对应于函子
\[
  h_V=\operatorname{Hom}_{\mathscr V}(-,V)
\]
的函子都与投射极限交换。因此，对于每个有限长度的 \(S\)-概形 \(T\)，有函子性同构
\[
\begin{aligned}
  \operatorname{Hom}_{\mathscr S}(T,X)
  &\simeq
  \varprojlim_i\operatorname{Hom}_{\mathscr S}(T,X_i)
\\
  &\simeq
  \varprojlim_i\operatorname{Hom}_{\mathscr V}(T,\widehat X_i)
  \simeq
  \operatorname{Hom}_{\mathscr V}
    \left(T,\varprojlim_i\widehat X_i\right).
\end{aligned}
\]
因此，函子 \(X\mapsto\widehat X/\widehat S\) 与
\((\operatorname{Sch}/S)\) 中存在的投射极限交换，特别地与有限投射极限交换。
\hfill\(\square\)

% 下一断点： Exp7B 原文件页码 26, 合并阅读版页码 539,
% 再版印刷页码 528, 紧接其前： 1.3.
